\section{Successive scales and relative spectral transport}
\label{sec:hierarchies}

The relative form estimates determine more than the
ordered small eigenvalues. They transport groups of
marked eigenspaces whenever those groups have a
relative gap. For fixed restriction maps, the
limiting groups have a direct description by
successive kernels.

\subsection{Successive kernels of fixed maps}

Let $V$ be a finite-dimensional real space with
positive mass $M$. For fixed linear maps $d_j$
to Euclidean spaces, put $B_j(x)=\|d_jx\|^2$
and
\[
 K_0=V,\quad K_j=\bigcap_{i\le j}\ker d_i,\quad
 V_j=K_{j-1}\cap K_j^{\perp_M},\quad
 r_j=\dim K_{j-1}-\dim K_j.
\]
Zero values of $r_j$ are allowed. Suppose
$t_j>0$, $t_{j+1}/t_j\to0$, and
$E_p=\sum_{j=1}^st_jB_j$. Let $Q_p,N_p$
satisfy
\begin{equation}\label{eq:hierarchy-relative-pair}
 (1-e_p)E_p\le Q_p\le(1+e_p)E_p,\qquad
 (1-e_p)M\le N_p\le(1+e_p)M,\qquad e_p\to0.
\end{equation}

\begin{theorem}\label{thm:fixed-map-hierarchy}
For each $j$ with $r_j>0$, precisely $r_j$
positive eigenvalues of $(Q_p,N_p)$ have
scale $t_j$. Dividing them by $t_j$ gives,
with multiplicity, the positive generalized
eigenvalues of
\begin{equation}\label{eq:hierarchy-leading-pair}
 (B_j|_{K_{j-1}},M|_{K_{j-1}}).
\end{equation}
The exact kernel has dimension $\dim K_s$.
The cumulative spectral space of all groups
with index greater than $j$, together with
the kernel, tends to $K_j$. The $j$-th
group tends to $V_j$ in the $M$-Grassmannian.
\end{theorem}
\begin{proof}
For fixed $j$ set
$S_{p,j}=(N_p+Q_p/t_j)^{-1}N_p$.
The vector $u=S_{p,j}f$ solves
\begin{equation}\label{eq:hierarchy-resolvent}
 N_p(u,x)+t_j^{-1}Q_p(u,x)=N_p(f,x).
\end{equation}
Testing with $x=u$ bounds $\|u\|_M$ and
$E_p(u)/t_j$ by constants times
$\|f\|_M$ and $\|f\|_M^2$, respectively.
For $j>1$, positivity of
$\sum_{i<j}B_i$ on $K_{j-1}^{\perp_M}$
then gives
\[
 \operatorname{dist}_M(u,K_{j-1})^2
       \le C(t_j/t_{j-1})\|f\|_M^2.
\]
For $j=1$ there is no excluded complement.

The form estimate implies its bilinear
version
\begin{equation}\label{eq:relative-bilinear-error}
 |(Q_p-E_p)(u,x)|
       \le e_p\sqrt{E_p(u)}\sqrt{E_p(x)}.
\end{equation}
This follows by passing to the quotient by
$\ker E_p$ and conjugating by its positive
square root. For $x\in K_{j-1}$, all
$B_i(u,x)$ with $i<j$ vanish exactly,
and $E_p(x)/t_j$ is bounded. Divide
\eqref{eq:hierarchy-resolvent} by the
scales as written and pass to a subsequence.
Every limit $u\in K_{j-1}$ satisfies
\[
 M(u,x)+B_j(u,x)=M(f,x),\qquad x\in K_{j-1}.
\]
This equation has a unique solution.
Finite dimension therefore gives operator
norm convergence of $S_{p,j}$ to the
$M$-orthogonal compression of
$(I+B_j|_{K_{j-1}})^{-1}$, extended by
zero on the complement. Conjugating by
the positive mass square roots turns this
into ordinary self-adjoint norm convergence.

The eigenvalues of $S_{p,j}$ are
$(1+\lambda/t_j)^{-1}$. Their limits are
zero on the fast complement, the $r_j$
values strictly between zero and one
specified by \eqref{eq:hierarchy-leading-pair},
and one on $K_j$. This proves the scale
and coefficient assertions. Repeating for
all $j$ accounts for every dimension.
The exact kernel is $K_s$ by
\eqref{eq:hierarchy-relative-pair}.

The isolated limit eigenvalue one gives
convergence of the cumulative slow projector
to the projector onto $K_j$. Differences
of consecutive cumulative projectors give
the projector onto
$K_{j-1}\cap K_j^{\perp_M}$. This proves
the subspace assertions, including steps
with $r_j=0$.
\end{proof}

The pair \eqref{eq:hierarchy-leading-pair}
can equivalently be taken on
$K_{j-1}/K_j$, with its quotient mass
$\inf_{k\in K_j}M(x+k)$.
If positive row weights satisfy
$w_e/t_j\to c_e\in(0,\infty)$ within
each group, use
$B_j=\sum_{e\text{ in group }j}c_ea_e^*a_e$.
The relative error of this positive sum
tends to zero without dividing by its
smallest weight. Thus logarithmic and
iterated-logarithmic scales are retained.
The grouping is attached to a sector
or path with those limits. It need not
define one filtration on the entire
multiparameter base.

\subsection{A relative-gap estimate}

At a single parameter, let $E\ge0$, $M>0$
on a space of dimension at most $n$, and
suppose $Q,N$ obey
\[
 (1-e)E\le Q\le(1+e)E,\qquad
 (1-e)M\le N\le(1+e)M,\qquad 0<e<1/4.
\]
Order their generalized eigenvalues
$\gamma_i$ and $\lambda_i$ in the same
direction, including zeros. A consecutive
positive index set $I$ is relatively
separated by $\kappa$ if $0<\kappa\le1$
and
\begin{equation}\label{eq:relative-gap}
 |\gamma_i-\gamma_j|
       \ge\kappa\max(\gamma_i,\gamma_j),
       \qquad j\in I,\ i\notin I.
\end{equation}
Let $F_I$ and $G_I$ be the corresponding
finite and observed subspaces.

\begin{lemma}\label{lem:relative-subspace-transport}
If $e\le c_n\kappa$, then
\begin{equation}\label{eq:relative-subspace-transport}
 \|P^M_{F_I}-P^M_{G_I}\|_M\le C_ne/\kappa.
\end{equation}
The exact kernels coincide. There is no
factor depending on the positive spectral
condition number.
\end{lemma}
\begin{proof}
Choose an $M$-orthonormal eigenbasis $u_i$
of $(E,M)$ and an $N$-orthonormal eigenbasis
$v_j$ of $(Q,N)$. Min--max gives
$\lambda_j/\gamma_j$ between
$(1-e)/(1+e)$ and $(1+e)/(1-e)$
for positive eigenvalues, with zeros
preserved. Testing the $v_j$ equation
against $u_i$ yields
\[
 (\gamma_i-\lambda_j)M(u_i,v_j)
   =-(Q-E)(u_i,v_j)+\lambda_j(N-M)(u_i,v_j).
\]
Since $E(v_j)\le\lambda_j/(1-e)$ and
$M(v_j)\le(1-e)^{-1}$, the right side
is at most
$Ce(\sqrt{\gamma_i\lambda_j}+\lambda_j)$
in absolute value. For $j\in I$, $i\notin I$,
\eqref{eq:relative-gap} and smallness of
$e/\kappa$ give
\[
 |\gamma_i-\lambda_j|
       \ge c\kappa\max(\gamma_i,\lambda_j).
\]
Every off-block coefficient is therefore
at most $Ce/\kappa$. A unit combination
of at most $n$ columns has off-block
norm at most $C_ne/\kappa$, after using
mass equivalence. Equal dimensions and
the principal-angle identity give
\eqref{eq:relative-subspace-transport}.
Explicitly, an $M$-orthonormal basis of
$G_I$ whose off-block norm is at most
$a$ has projection to $F_I$ with least
singular value at least $\sqrt{1-a^2}$;
its largest principal-angle sine is at
most $a$. Kernel equality follows
directly from the two-sided energy
inequality.
\end{proof}

Through \eqref{eq:actual-fixed-observer},
the lemma compares marked actual subspaces
with the independently constructed finite
subspaces. It does not identify forms
on unrelated manifolds without that map.
Individual vectors require an internal
relative gap; repeated leading eigenvalues
give subspace assertions.

\subsection{Limits of the fixed-map description}

The fixed-map hypothesis has mathematical
content. With identity mass and target
weight, the moving map
\[
 d_h=\begin{pmatrix}1&0\\1&h^m\end{pmatrix},
 \qquad m>0,
\]
has row norms bounded away from zero,
but $d_h^*d_h$ has trace $2+h^{2m}$
and determinant $h^{2m}$. Its eigenvalues
are asymptotic to $2$ and $h^{2m}/2$.
Multiplying the energy by $t(h)\to0$
gives two small eigenvalue scales $2t$ and
$th^{2m}/2$. The rate at which the
rows become dependent is additional
data. Replacing the lower-right entry
by $h^m-h^m+h^n$, $n>m$, replaces the
second exponent by $2n$; cancellation
must be performed before assigning
orders to minors.

The mass is equally necessary. For
$E_h=h^2I$ and $M_h=\operatorname{diag}(h^2,1)$,
the generalized eigenvalues are $1,h^2$.
Vanishing energy on both coordinate
vectors does not make both eigenvalues tend to zero.
Finally, orthogonal conjugates of
$\operatorname{diag}(h,h^2)$ have identical
characteristic coefficients, but a
nonconvergent rotation angle produces
nonconvergent lines. Eigenvalues alone
do not determine marked spectral
subspaces.

For the conifold theorem the restriction
rows are fixed and the central mass is
positive definite. Thus the hypotheses
of Theorem~\ref{thm:fixed-map-hierarchy}
are verified on every capacity-ratio
sector. The two-dimensional pair also
has a relative gap throughout the full
domain, which gives a stronger pointwise
eigenline statement.
